\chapter{Human Nature as a Managed Narrative}
\label{ch:human-nature}

\epigraph{It is not that people are either good or bad by nature, but that they can be made more selfish or more cooperative by the stories they are told about themselves.}{Samuel Bowles}

\noindent
This is the signature case study of the book.
The belief that ``human nature is selfish'' is among the most consequential doxonomic products of the modern era: an anthropological regime reinforced through centuries of institutional investment, whose self-fulfilling properties make it appear to confirm itself.
Whether it is \emph{the} most consequential---more than racial hierarchy, national myth, or religious cosmology---is a comparative question this chapter does not attempt to settle.
What it does attempt is a detailed doxonomic analysis of one regime whose mechanisms are well-documented and whose effects are measurable.

The analysis proceeds in stages: first the empirical evidence on what human nature actually looks like (Section~\ref{sec:ch21-evidence}); then the institutional genealogy of the selfishness regime (Sections~\ref{sec:ch21-genealogy}--\ref{sec:ch21-media}); the self-fulfilling mechanism by which the belief generates its own confirmation (Section~\ref{sec:ch21-self-fulfilling}); the policy feedback loop (Section~\ref{sec:ch21-policy}); the formal competition model (Section~\ref{sec:ch21-competition}); counter-narratives and their institutional base (Section~\ref{sec:ch21-counter}); and cross-cultural comparison (Section~\ref{sec:ch21-cross-cultural}).

\section{The Claim and the Evidence}
\label{sec:ch21-evidence}

\subsection{The dominant anthropological regime}

The dominant anthropological regime in Western market societies holds that:
\begin{enumerate}[nosep]
  \item humans are fundamentally self-interested,
  \item cooperation requires external enforcement (contracts, laws, incentives),
  \item altruism is either disguised self-interest or evolutionarily marginal,
  \item competition is the natural state; solidarity must be artificially maintained.
\end{enumerate}

This regime is not presented as ideology; it is presented as science---as the sober conclusion of economics, evolutionary biology, and rational choice theory.
Its power lies precisely in this scientific self-presentation: to question it is to be naive, sentimental, or ignorant of ``how the world really works.''

\subsection{What the experimental evidence actually shows}

The experimental evidence does not support these claims as a complete anthropology, though it does not deny that self-interest, status competition, and opportunism are real human tendencies.

\paragraph{Public goods games.}
\textcite{fischbacher2001} showed that the majority of participants in public goods games are \emph{conditional cooperators}: they cooperate when they expect others to cooperate.
Only about 30\% behave as consistent free riders.
The remaining 20\% display mixed or erratic patterns.
This means the modal human is not selfish---the modal human is \emph{conditionally pro-social}, adjusting behavior to perceived social norms.

\paragraph{Cross-cultural evidence.}
\textcite{henrich2001} demonstrated generosity in ultimatum and dictator games across 15 small-scale societies on four continents.
People in diverse societies routinely reject pure self-interest, though the degree varies with market integration and institutional context.
No society tested exhibited the pure self-interest predicted by the homo economicus model.

\paragraph{Commons management.}
\textcite{ostrom1990} documented hundreds of communities successfully managing common-pool resources---fisheries, irrigation systems, forests, grazing lands---without privatization or state enforcement, many for centuries.
The ``tragedy of the commons'' \autocite{hardin1968} turns out to be not a law of nature but a prediction about behavior under specific institutional conditions (open access with no communication), conditions that real communities routinely avoid through institutional design.

\paragraph{Developmental evidence.}
Children as young as 14 months display helping behavior (picking up dropped objects for an adult) without training or reward.
By age 3, children enforce fairness norms on others and protest unfair distributions.
By age 5, they engage in costly punishment of cheaters.
These behaviors appear before cultural training has had much time to operate, suggesting an evolved pro-social foundation.

\paragraph{Neuroscience.}
Cooperation activates reward centers in the brain (striatum, ventromedial prefrontal cortex) even when the cooperator receives no material benefit.
Rejecting unfair offers in ultimatum games activates the anterior insula, a region associated with disgust.
The neural architecture supports cooperation as intrinsically rewarding, not merely instrumentally advantageous.

\paragraph{The synthetic picture.}
\textcite{bowles2011} synthesized this evidence: cooperation is a fundamental, evolved human capacity, coexisting with competition and conflict rather than replacing them.
Humans are \emph{conditional cooperators with a taste for punishing cheaters}---a combination that supports cooperation in groups but also generates inter-group conflict.
The best available anthropology is not ``humans are selfish,'' not ``humans are cooperative,'' but \emph{humans are flexibly social}, with the balance between cooperation and competition depending heavily on institutional design, narrative environment, and perceived norms \autocite{gintis2005}.

The doxonomic point is not that the selfishness regime is ``wrong'' in every particular---self-interest is real, free-riding occurs, enforcement helps.
The point is that the regime \emph{selectively amplifies} one real tendency while suppressing evidence of others, producing a distorted anthropology that serves specific institutional interests.

\section{The Institutional Genealogy of Homo Economicus}
\label{sec:ch21-genealogy}

\subsection{Intellectual precursors (1600--1900)}

\begin{definition}[Homo Economicus as Anthropological Regime]
\label{def:homo-economicus}
\index{homo economicus}
\emph{Homo economicus} is the anthropological regime $\alpha_{\text{HE}}$ defined by:
motivational theory = rational self-interest;
social ontology = isolated individuals;
capacity claim = collective action requires external incentives.
\end{definition}

The intellectual lineage:
\begin{itemize}[nosep]
  \item \textbf{Hobbes} (1651): ``the life of man, solitary, poor, nasty, brutish, and short''---without a sovereign, self-interest produces war.
  \item \textbf{Mandeville} (1714): \emph{The Fable of the Bees}---``private vices, public benefits.'' Self-interest, properly channeled, produces social good.
  \item \textbf{Smith} (1776): ``It is not from the benevolence of the butcher, the brewer, or the baker that we expect our dinner, but from their regard to their own interest.'' Note: Smith also wrote \emph{The Theory of Moral Sentiments} (1759), emphasizing sympathy and social connection, but this text is far less cited in economics curricula.
  \item \textbf{Bentham} (1789): utility maximization as the foundation of rational choice.
  \item \textbf{Spencer} (1864): ``survival of the fittest'' applied to human societies---Social Darwinism as a naturalization of competition.
\end{itemize}

The doxonomic observation: none of these thinkers simply ``discovered'' that humans are selfish.
Each was writing in a specific institutional context---the rise of market society, the dissolution of feudal obligations, colonial expansion---and their theories served specific functions in legitimating the emerging order.

\subsection{Formalization (1900--1960)}

The selfishness assumption was transformed from a philosophical claim into a formal axiom:
\begin{itemize}[nosep]
  \item \textbf{Pareto} (1906): ordinal utility theory---agents have well-ordered preferences, interpreted as self-regarding.
  \item \textbf{Robbins} (1932): economics as the science of allocating scarce means among competing ends---no room for altruism, solidarity, or non-instrumental motivation.
  \item \textbf{Von Neumann and Morgenstern} (1944): expected utility theory---the rational agent maximizes expected payoff.
  \item \textbf{Arrow and Debreu} (1954): general equilibrium---the economy as a system of self-interested agents reaching optimal allocation through markets.
  \item \textbf{Friedman} (1953): the ``as if'' defense---even if people are not self-interested, models that assume they are predict well, so the assumption is ``vindicated.''
\end{itemize}

The formalization was crucial doxonomically: it turned a contestable philosophical claim (``people are selfish'') into a mathematical axiom (``preferences are complete and transitive'') that appeared value-neutral.
Students who might have questioned the philosophy were instead taught the mathematics, absorbing the anthropology as an unstated background assumption.

\subsection{Institutional build-out (1960--2000)}

The second half of the 20th century saw the selfishness regime embedded in an institutional infrastructure of unprecedented scale:

\begin{enumerate}
  \item \textbf{Economics departments.}
  From roughly 200 PhD-granting economics departments in the US alone, graduating $\sim$1,200 PhDs per year, who then teach the self-interest axiom to $\sim$1.5 million undergraduate students annually.

  \item \textbf{Business schools.}
  The MBA program---which reached $\sim$200,000 graduates per year in the US by 2000---teaches principal-agent theory, incentive design, and shareholder value maximization, all premised on self-interest.
  Jensen and Meckling's (1976) ``Theory of the Firm'' paper, which models managers as self-interested agents who must be disciplined by contracts, became the most cited paper in economics.

  \item \textbf{Think tanks.}
  Free-market think tanks (Heritage, Cato, AEI, IEA, Adam Smith Institute, Atlas Network of 500+ affiliated organizations) produce a steady stream of publications, op-eds, and policy recommendations premised on the rational self-interest model.

  \item \textbf{Textbooks.}
  Mankiw's \emph{Principles of Economics}---the dominant introductory textbook from the mid-1990s onward---presents self-interest as Principle \#4 (``People respond to incentives'') and frames market outcomes as the natural result of rational choice.
  Over 4 million copies sold.
\end{enumerate}

\begin{example}
\label{ex:funding-estimate}
A rough doxonomic accounting of the selfishness-regime's annual institutional investment in the United States (circa 2010):

\begin{center}
\begin{tabular}{lrr}
\toprule
Channel & Annual expenditure (est.) & Annual reach (est.) \\
\midrule
Economics departments (teaching) & \$2--4 billion & 1.5 million students \\
Business schools (MBA + exec ed) & \$5--8 billion & 250,000 graduates \\
Free-market think tanks (top 50) & \$0.5--1 billion & Millions (media, policy) \\
Management consulting & \$2--3 billion & 100,000+ corporate clients \\
Business media (CNBC, Bloomberg, etc.) & \$1--2 billion & 50+ million viewers \\
\midrule
\textbf{Total (conservative)} & \textbf{\$10--18 billion} & \\
\bottomrule
\end{tabular}
\end{center}

These figures are rough order-of-magnitude estimates, not precise accounting.
They include only the costs of institutions that teach or assume self-interest as a behavioral foundation; they exclude indirect channels (entertainment media, social media, interpersonal transmission).
Even at the lower bound, the annual investment in reproducing the selfishness regime exceeds the entire budget of the National Science Foundation (\$8.3 billion in 2020).
\end{example}

\section{The Education Channel}
\label{sec:ch21-education}

\subsection{Economics as doxonomic delivery mechanism}

The economics classroom is the single most important delivery mechanism for the selfishness regime.
It reaches more people, with more institutional authority, over a longer period of exposure, than any think tank or media outlet.

\begin{model}[Educational Doxonomic Pipeline]
\label{mod:education-pipeline}
\index{education pipeline}
The pipeline has three stages:
\begin{enumerate}[nosep]
  \item \textbf{Undergraduate economics} ($\sim$1.5M students/year in the US): Intro courses teach self-interest as the motivational axiom. Most students take only the intro course, absorbing the axiom without encountering the behavioral qualifications taught in advanced courses.
  \item \textbf{MBA programs} ($\sim$200K graduates/year): Principal-agent theory, incentive design, and shareholder primacy deepen the self-interest assumption and apply it to organizational design.
  \item \textbf{Corporate practice}: MBA graduates design compensation systems, performance metrics, and organizational structures premised on worker self-interest, creating institutional environments that elicit the behavior the model predicts.
\end{enumerate}
Each stage amplifies the previous one: the undergraduate absorbs the axiom, the MBA applies it, the executive institutionalizes it.
\end{model}

\subsection{The Frank-Gilovich-Regan finding}

\textcite{frank1993} documented that economics students behave more self-interestedly than students in other disciplines---contributing less in public goods games, defecting more in prisoner's dilemma, and expecting others to defect more.
The finding has been replicated across multiple countries and experimental protocols.

The causal question is contested (selection vs.\ indoctrination), but longitudinal evidence suggests both mechanisms operate.
Students who enter economics are somewhat more self-interested than average (selection), but their self-interested behavior increases during the course (indoctrination).
The indoctrination effect is concentrated in the first course, suggesting that the introductory axiom has a larger impact than advanced training.

\begin{proposition}[Economics Education as Doxonomic Treatment]
\label{prop:econ-education}
\index{economics education}
If the average introductory economics course shifts cooperative behavior by $\delta = -3$ to $-5$ percentage points (the range suggested by existing studies), and $\sim$1.5 million US students take the course annually, then the annual doxonomic output is:
\[
  \Delta_{\text{annual}} = 1.5 \times 10^6 \times \delta \approx 4.5\text{--}7.5 \times 10^6 \text{ cooperation-points lost per year.}
\]
Over a 40-year working life, the cumulative effect on the population stock of cooperative disposition is substantial---even before accounting for the multiplier effects of changed behavior influencing peers.
\end{proposition}

\subsection{What textbooks teach vs.\ what the evidence says}

\begin{center}
\begin{tabular}{p{5.5cm}p{5.5cm}}
\toprule
Standard textbook claim & Current evidence \\
\midrule
People are rational self-maximizers & Most people are conditional cooperators; pure maximizers are a minority \\[4pt]
Altruism is anomalous & Pro-social behavior is developmentally early, neurologically grounded, cross-culturally robust \\[4pt]
Commons are tragic without privatization & Hundreds of self-governing commons have persisted for centuries \\[4pt]
Incentives always help & Monetary incentives can crowd out intrinsic motivation and reduce cooperation \\[4pt]
Markets produce optimal outcomes & Markets produce efficient outcomes only under restrictive conditions rarely met in practice \\
\bottomrule
\end{tabular}
\end{center}

The gap between textbook claims and current evidence is not a gap that exists because the evidence is new.
Ostrom published in 1990.
Fischbacher and colleagues published in 2001.
Henrich and colleagues published in 2001.
These findings are a generation old.
Their absence from standard introductory textbooks is itself a doxonomic phenomenon---the institutional inertia of a curriculum controlled by departments whose senior faculty were trained in the self-interest paradigm.

\section{Media and Popular Culture}
\label{sec:ch21-media}

The selfishness regime extends far beyond the classroom.
Popular culture provides continuous reinforcement through narrative channels that reach audiences who never take an economics course.

\paragraph{Financial media.}
CNBC, Bloomberg, the Financial Times, and the Wall Street Journal operate within the self-interest framework as default.
Market movements are explained by ``rational'' responses to incentives; cooperation is a ``market failure''; regulation is ``interference'' with natural market processes.
The aggregate audience of financial media in the US exceeds 50 million.

\paragraph{Entertainment.}
The ``greed is good'' speech in \emph{Wall Street} (1987) became cultural shorthand for a worldview already institutionalized.
Reality television (e.g., \emph{Survivor}, \emph{The Apprentice}) is structured around competition, betrayal, and individual victory---reinforcing the narrative that life is a zero-sum game.
Self-help literature frequently adopts evolutionary-psychology framing: ``your brain is wired for competition.''

\paragraph{Evolutionary psychology popularization.}
Books like Dawkins' \emph{The Selfish Gene} (1976)---a brilliant work of science communication whose title was widely misread as a claim about organisms rather than genes---provided scientific-sounding support for the selfishness regime.
The ``gene's-eye view'' was translated, in popular understanding, into ``selfishness is natural, cooperation is an illusion.''
Dawkins himself has protested this misreading, but the doxonomic damage was done: the phrase ``selfish gene'' became a cultural reference point for the naturalization of self-interest.

\paragraph{The quantitative asymmetry.}
Narratives about human selfishness and competition are vastly more prevalent in mainstream media than narratives about human cooperation and solidarity.
A content analysis of US newspaper opinion pages would reveal that ``incentive'' appears orders of magnitude more frequently than ``solidarity,'' and ``market'' more frequently than ``commons.''
This asymmetry is not a natural reflection of reality; it reflects the institutional infrastructure that produces media content.

\section{The Self-Fulfilling Prophecy}
\label{sec:ch21-self-fulfilling}

\subsection{The mechanism}

Consider a first-year university student taking an introductory economics course.
The professor explains that rational agents maximize their own utility; the textbook models assume self-interest; the problem sets reward strategic thinking.
The student leaves the course believing---not because she has been argued into it, but because it is the unquestioned premise of every exercise she completed---that most people are selfish.
She enters a group project expecting free-riders.
She monitors her partners for shirking.
Her vigilance makes them uncomfortable; trust erodes; cooperation falters.
``See?'' she concludes.
``People really are selfish.''
She has just experienced the most dangerous property of the selfishness regime: it generates the evidence for its own truth.

\begin{proposition}[Doxonomic Self-Fulfillment]\footnote{This result follows from the model assumptions; it is a heuristic prediction, not an empirically validated law. The qualitative mechanism (narrative shifts norm perceptions, which shifts behavior, which confirms the narrative) is well-supported; the specific dynamics are stylized.}
\label{thm:self-fulfilling}
\index{self-fulfilling prophecy}
Let agents be conditional cooperators who cooperate when $\hat{\mu}_j(\text{cooperate}) > \tau_j$, where $\hat{\mu}_j$ is $j$'s perception of the cooperation norm and $\tau_j$ is $j$'s cooperation threshold.
If a doxonomic intervention reduces $\hat{\mu}_j(\text{cooperate})$ by shifting perceived norms (by amount $\delta$), then:
\begin{enumerate}[nosep]
  \item fewer agents exceed their threshold: cooperation drops,
  \item the observed drop confirms the narrative (``see, people are selfish''),
  \item which further reduces $\hat{\mu}_j$ for those who observe behavior,
  \item producing a downward cascade until the system reaches a new (lower) equilibrium.
\end{enumerate}
Formally, the dynamics satisfy:
\[
  \mu^{(t+1)}_{\text{coop}} = G\!\left(\hat{\mu}^{(t)}_{\text{coop}} - \delta\right)
\]
where $G$ is the fraction of the population whose threshold $\tau_j$ falls below the perceived norm.
If $G'(x) > 1$ at the initial equilibrium, the doxonomic shock $\delta$ triggers a cascade to the low-cooperation equilibrium.
\end{proposition}

\subsection{Multiple equilibria}

The self-fulfilling mechanism implies that the same population, with the same underlying preferences, can sustain \emph{different equilibrium cooperation levels} depending on the narrative environment:

\begin{itemize}[nosep]
  \item \textbf{High-cooperation equilibrium}: most people believe others cooperate $\to$ most cooperate $\to$ belief confirmed.
  \item \textbf{Low-cooperation equilibrium}: most people believe others defect $\to$ most defect $\to$ belief confirmed.
\end{itemize}

Both equilibria are self-consistent.
The narrative environment---specifically, which stories about human nature are institutionally amplified---determines which equilibrium the population settles into.
The selfishness regime pushes toward the low-cooperation equilibrium; a cooperative regime would push toward the high-cooperation equilibrium.

\begin{example}
\label{ex:self-fulfilling-numerical}
Let $\tau_j \sim \text{Beta}(3, 2)$ (most people have moderate cooperation thresholds, with a mean of 0.6).
At the high-cooperation equilibrium, $\hat{\mu} = 0.75$: 86\% of agents have $\tau_j < 0.75$ and cooperate.
At the low-cooperation equilibrium, $\hat{\mu} = 0.30$: only 17\% cooperate.

A doxonomic shock of $\delta = 0.15$ (shifting perceived norms from 0.75 to 0.60) triggers a cascade: at $\hat{\mu} = 0.60$, 66\% cooperate.
But the observed 66\% further reduces $\hat{\mu}$, which reduces cooperation further, until the system settles at the low-cooperation equilibrium ($\hat{\mu} \approx 0.30$, 17\% cooperate).

The \emph{same population} has gone from 86\% to 17\% cooperation---not because preferences changed, but because the narrative environment tipped the norm perception past the critical threshold.
\end{example}

\subsection{The doxonomic trap}

\begin{definition}[Doxonomic Trap]
\label{def:doxonomic-trap}
\index{doxonomic trap}
A \emph{doxonomic trap} is a locally stable equilibrium that is Pareto-inferior to an alternative equilibrium but is sustained by the self-fulfilling properties of the dominant narrative.
Escaping the trap requires a coordinated shift in both narrative and norm perception---which is difficult precisely because the current narrative makes the alternative seem naive or impossible.
\end{definition}

The selfishness regime constitutes a doxonomic trap: the low-cooperation equilibrium is Pareto-inferior (everyone would be better off with more cooperation), but the narrative that ``people are selfish'' is self-confirming, and anyone who proposes the cooperative alternative is dismissed as ignoring ``human nature.''

\section{Affect, Precision, and the Moral Basin}
\label{sec:ch21-affect-precision}

The self-fulfilling mechanism of Section~\ref{sec:ch21-self-fulfilling} operates through perceived norms: the regime shifts what agents expect others to do, and behavior follows expectation.
A deeper layer operates through moral feeling itself.
The selfishness regime does not only tell agents that others will defect; it also teaches them which emotional responses count as \emph{moral warrant}.
Contempt for a free-rider reads as moral perception; unease at a cooperator who ``left money on the table'' reads as prudence; sympathy for a diffuse victim reads as sentimentality.
Over time the regime shapes not only propositional belief but the precision assigned to affect---which gut responses are treated as evidence, and which are treated as noise to be overridden.

This section sketches an individual-level mechanism that feeds the population dynamics already developed.
Doxonomics is not psychology (Chapter~\ref{ch:ontology-belief}), but population belief aggregates individual states, and some of the most durable regimes operate by recruiting the machinery of moral affect.
The account is stylized and active-inference-compatible; it is a proposed synthesis, not a canonical result.

\subsection{Feelings as dashboards, not verdicts}

The folk theory of moral cognition treats feeling as perception: disgust reveals wrongness, anger reveals injustice, guilt reveals culpability.
The empirical picture is more constrained.
Moral evaluation can, in some tasks, be reported faster than self-reported felt emotion, and incidental affect (disgust primed by a bad smell, anger carried over from traffic) leaks into unrelated moral judgments, though meta-analysis suggests average incidental-affect effects are small.
The useful compression is that moral feelings are \emph{dashboards}---fast compressions of harm, threat, norm violation, and action urgency---not \emph{verdicts}.

Under active inference---a family of frameworks in which the brain is modeled as continually predicting its sensory inputs and updating to minimize prediction error---emotion is an \emph{interoceptive inference}: the organism's best model of its own bodily state (heart rate, arousal, visceral tone), used to regulate action.
What governs how much that inference shapes downstream judgment is \emph{precision}---the confidence the system assigns to the bodily signal relative to other evidence channels.
When precision on interoception is high and evidence on harm, intent, or victim structure is diffuse, the system falls into a moral basin driven by feeling.
When precision is more evenly distributed, the same feeling becomes one input among several.

\subsection{A minimal state space}

Let the individual moral state be
\[
  x_t = \bigl[\,I,\ \Pi,\ H,\ T,\ N,\ V,\ C,\ R,\ S,\ U\,\bigr]
\]
with $I$ interoceptive arousal, $\Pi$ the precision placed on it, $H$ perceived harm, $T$ inferred intent, $N$ taboo/purity salience, $V$ victim salience, $C$ self-causation, $R$ repair efficacy, $S$ social confirmation, and $U$ residual ambiguity.
The state moves on a potential $\mathcal{V}(x; c)$ shaped by context $c$: threat cues, narrative framing, group signals, institutional environment.
The picture is a landscape: $\mathcal{V}$ has local minima---\emph{attractors}---around which the system settles, each surrounded by a \emph{basin} of nearby states that drain into it.
Attractors here are not physical locations but stable patterns of appraisal and disposition.
Candidate attractors correspond to action dispositions: \emph{avoid} (disgust-like exclusion), \emph{punish} (anger-like sanction), \emph{repair} (guilt-like restitution), \emph{help} (empathic concern), and \emph{suspend} (withholding judgment).
Context $c$ reshapes the landscape: a regime that deepens one basin and flattens others changes which attractor the same facts tend to produce.
The same person, facing the same facts, can land in different basins depending on which variables the environment amplifies.

\begin{definition}[Moral Distortion Index]
\label{def:moral-distortion}
\index{moral distortion index}
\index{distortion index, moral}
For an individual in moral state $x$, the \emph{moral distortion index} is
\[
  D(x) \;=\; \frac{\Pi\,I + S}{H + T + V + C + R + \varepsilon}
\]
for small $\varepsilon > 0$.
$D \ll 1$: judgment is anchored to appraisal evidence (harm, intent, victim structure, causation, repair).
$D \gg 1$: bodily affect weighted by precision, plus social reinforcement, outruns appraisal.
\end{definition}

$D$ is a bookkeeping device, not a measurement instrument.
It names the regime in which ``this feels wrong'' most readily hardens into ``this \emph{is} wrong,'' and---symmetrically---the regime in which moral feeling serves as useful input rather than mistaken revelation.

\subsection{Doxonomic amplification of \texorpdfstring{$D$}{D}}

A regime does not control the full state $x$; it shapes context $c$, and context shapes which variables the environment makes salient.
The selfishness regime raises $D$ asymmetrically through two complementary moves:
\begin{enumerate}[nosep]
  \item \textbf{Amplification.} It elevates $\Pi\,I$ and $S$ on regime-supporting stimuli---betrayal narratives about free-riders, affective framings of welfare recipients, spectacle-driven coverage of individual wrongdoing---so that anger and contempt feel like moral perception.
  \item \textbf{Diffusion.} It suppresses $H$, $V$, $C$, and $R$ on regime-threatening stimuli---institutional harms spread across populations, victims with no camera, causal chains too long to narrate, reparative policies framed as fiscally irresponsible.
\end{enumerate}
The aggregate effect is to grant moral feeling high precision where it supports the regime and low precision where it would challenge it.
This composes with Section~\ref{sec:ch21-self-fulfilling}: narrative shifts norm perception, and it also shifts which feelings are granted moral authority.

\begin{proposition}[Affective Recruitment]\footnote{This proposition follows from the model assumptions; it is a structural claim about how regimes exploit affect, not a calibrated prediction.}
\label{prop:affective-recruitment}
\index{affective recruitment}
A regime that selectively amplifies $\Pi\,I$ and $S$ on regime-supporting stimuli while diffusing $H$, $V$, $C$, and $R$ on regime-threatening stimuli raises population $D$ asymmetrically.
The resulting distribution of moral feeling tracks institutional interest rather than evidential warrant, even when no agent experiences the shift as manipulation.
\end{proposition}

The mechanism matters precisely because it is not experienced as indoctrination.
Agents feel their moral responses as their own.
The regime has shaped not what they think but what feels morally self-evident---which is a deeper lock-in than propositional belief, and correspondingly harder to dislodge.

\subsection{Hysteresis and the cost of exit}

Attractor dynamics imply hysteresis: once a moral basin deepens through repetition, social reinforcement, and enacted judgment, the evidence required to exit it exceeds the evidence that would have been required to prevent entry.
This is consistent with the lock-in properties of doxonomic systems developed in Chapter~\ref{ch:formal-models}.
The policy consequence is that counter-doxonomic work cannot rely on argument alone.
Arguments operate on appraisal variables ($H$, $T$, $V$); lock-in operates on the affect-precision channel $\Pi\,I$ and the social-coupling channel $S$.
Shifting the moral basin requires sustained experiential counter-evidence (cooperative institutions that visibly work), sustained reframing of affect (narratives that re-locate anger toward structural sources of harm and concern toward structurally produced victims), or both.

Whether populations actually exhibit the asymmetric $D$ predicted by Proposition~\ref{prop:affective-recruitment} is an open empirical problem.
Candidate approaches are discussed in Chapter~\ref{ch:doxometric-measurement} (individual-level affect measurement) and Chapter~\ref{ch:text-media} (regime-level amplification patterns in discourse).

\section{The Policy Feedback Loop}
\label{sec:ch21-policy}

The selfishness regime does not merely describe behavior; it shapes the institutional environment in which behavior occurs, creating a feedback loop:

\begin{model}[Policy Feedback]
\label{mod:policy-feedback}
\index{policy feedback}
\begin{enumerate}[nosep]
  \item The selfishness assumption enters policy design: regulators assume that agents will cheat unless monitored; managers assume workers will shirk unless incentivized; policymakers assume citizens will free-ride unless forced to contribute.
  \item Policy instruments are designed accordingly: surveillance systems, performance-based pay, privatization, means-testing.
  \item These instruments signal to citizens that they are not trusted, which reduces intrinsic motivation and cooperative behavior \autocite{frey1997}.
  \item The resulting behavior confirms the selfishness assumption, justifying further monitoring and incentivization.
\end{enumerate}
\end{model}

\textcite{frey1997} documented the ``crowding out'' effect: monetary incentives can reduce intrinsic motivation.
Offering to pay people for blood donations \emph{reduces} donations compared to asking them to donate voluntarily.
Imposing fines for late pickup at daycare centers \emph{increases} lateness (parents treat the fine as a price, replacing a social norm with a market transaction).

The policy feedback loop means that the selfishness regime is not just a belief about behavior---it is a \emph{generator} of behavior.
The regime creates the institutional conditions under which self-interest is the rational response, then points to the resulting behavior as evidence that self-interest is natural.

\section{Two Competing Regimes: A Formal Model}
\label{sec:ch21-competition}

\begin{model}[Anthropological Regime Competition]
\label{mod:regime-competition}
\index{anthropological regime!competition}
Let $\belief{S}(t)$ be the prevalence of the selfish regime and $\belief{C}(t) = 1 - \belief{S}(t)$ the cooperative alternative.
The dynamics are:
\begin{align}
  \frac{d\belief{S}}{dt} &= \belief{S}\bigl(\alpha_S \phi_S - \gamma_S - \beta_{SC}\belief{C}\bigr) \\
  \frac{d\belief{C}}{dt} &= \belief{C}\bigl(\alpha_C \phi_C - \gamma_C - \beta_{CS}\belief{S}\bigr)
\end{align}
where:
\begin{itemize}[nosep]
  \item $\phi_S, \phi_C$ are institutional funding levels ($\phi_S \gg \phi_C$),
  \item $\alpha_S, \alpha_C$ are conversion efficiencies (funding $\to$ belief change),
  \item $\gamma_S, \gamma_C$ are natural decay rates,
  \item $\beta_{SC}, \beta_{CS}$ are competitive interaction coefficients.
\end{itemize}
\end{model}

\begin{proposition}[Funding Threshold for Regime Change]
\label{prop:funding-threshold}
\index{regime change!funding threshold}
The cooperative regime displaces the selfish regime if and only if
\[
  \alpha_C \phi_C > \gamma_C + \beta_{CS} \cdot \belief{S}^*
\]
where $\belief{S}^*$ is the current selfish-regime prevalence.
Given the current dominance of the selfish regime ($\belief{S}^* \approx 0.7$), this implies:
\[
  \phi_C > \frac{\gamma_C + 0.7 \beta_{CS}}{\alpha_C}.
\]
With plausible parameter values ($\gamma_C = 0.05$, $\beta_{CS} = 0.3$, $\alpha_C = 0.001$ per million dollars), the required annual investment is:
\[
  \phi_C > \frac{0.05 + 0.21}{0.001} = \$260 \text{ million/year.}
\]
This is a fraction of the selfishness regime's estimated \$10--18 billion annual investment, suggesting that the cooperative regime is not disadvantaged by intrinsic weakness but by massive resource asymmetry.
\end{proposition}

Note: these parameter values are illustrative, not empirically calibrated.
Rigorous estimation would require the methods of Chapters~\ref{ch:doxometric-measurement}--\ref{ch:causal-inference}.
The point is not the specific number but the \emph{structure} of the analysis: regime change depends on the ratio of institutional investment to natural decay and competitive pressure, and the required investment can in principle be estimated.

\section{Counter-Narratives and Their Institutional Base}
\label{sec:ch21-counter}

The selfishness regime has not gone unchallenged.
Several counter-narratives have emerged, each with its own (much smaller) institutional base:

\paragraph{Behavioral economics.}
Beginning with Kahneman and Tversky in the 1970s and gaining institutional power through Nobel Prizes (2002, 2017), the Journal of Behavioral Economics, and ``nudge units'' in government, behavioral economics challenges the rationality assumption---though it generally preserves self-interest, replacing ``rational selfishness'' with ``biased selfishness'' \autocite{kahneman2011}.

\paragraph{Commons research.}
Ostrom's work, recognized with the 2009 Nobel Prize, demonstrated that communities can manage shared resources without privatization.
The institutional base: the Ostrom Workshop (Indiana University), the International Association for the Study of the Commons, and a network of field researchers.
Funding level: tiny compared to the economics mainstream.

\paragraph{Cooperative economics.}
The cooperative movement (Mondragon, credit unions, worker cooperatives, platform cooperatives) provides empirical counter-evidence: organizations that assume cooperative motivation can be economically successful.
The institutional base is fragmented and underfunded compared to the corporate sector.

\paragraph{The funding asymmetry.}
A rough estimate of the counter-narrative institutional investment:

\begin{center}
\begin{tabular}{lr}
\toprule
Counter-narrative institution & Annual budget (est.) \\
\midrule
Behavioral economics research centers (top 10) & \$20--40 million \\
Commons research network (global) & \$5--10 million \\
Cooperative development organizations (US) & \$10--20 million \\
Progressive economic think tanks (EPI, Roosevelt, etc.) & \$30--50 million \\
\midrule
\textbf{Total counter-narrative investment} & \textbf{\$65--120 million} \\
\bottomrule
\end{tabular}
\end{center}

The funding ratio is approximately 100:1 in favor of the selfishness regime.
This asymmetry, not the relative quality of the evidence, is the primary explanation for the selfishness regime's dominance.

\section{Cross-Cultural Comparison}
\label{sec:ch21-cross-cultural}

The selfishness regime is not universal.
Different societies operate under different anthropological regimes, with different institutional infrastructures supporting different beliefs about human nature.

\paragraph{Japan.}
Japanese management theory emphasizes group harmony (\emph{wa}), lifetime employment, and organizational loyalty---an anthropological regime closer to ``humans are group-oriented'' than ``humans are selfish.''
The institutional base: the Ministry of Economy, Trade and Industry (METI), corporate culture, educational system emphasizing group work.
Note: this regime has been under pressure from neoliberal reform since the 1990s, illustrating the global reach of the selfishness regime's institutional infrastructure.

\paragraph{Nordic countries.}
Scandinavian societies maintain high levels of social trust and cooperation, supported by: universal welfare states, strong labor unions, cooperative traditions, and educational systems that emphasize solidarity.
The doxonomic infrastructure supporting cooperation is embedded in state institutions rather than private foundations.

\paragraph{Indigenous commons regimes.}
Many indigenous societies maintain anthropological regimes centered on reciprocity, kinship obligation, and commons stewardship.
These regimes have been systematically disrupted by colonial and neoliberal institutional intervention (cf.\ Chapter~\ref{ch:race-coloniality-gender}).

The cross-cultural comparison supports the doxonomic thesis: where institutional infrastructure supports cooperation, people cooperate more; where it supports self-interest, people behave more self-interestedly.
Human nature is not fixed at ``selfish'' or ``cooperative''---it is \emph{institutionally mediated}, and the mediating institutions are themselves products of material investment.

\section{Notes and References}
\label{sec:ch21-notes}

This chapter draws heavily on \textcite{bowles2011}, \textcite{fischbacher2001}, \textcite{henrich2001}, and \textcite{ostrom1990}.
On the ``tragedy of the commons'' and its limitations, see \textcite{hardin1968} and Ostrom's (1990) devastating response.
The developmental evidence for pro-social behavior is reviewed in Tomasello (2009).
Neuroscience of cooperation is surveyed in Rilling et al.\ (2002).

On the indoctrination effect of economics education, see \textcite{frank1993} and \textcite{bauman1996}.
The institutional history of homo economicus is traced in \textcite{mirowski2009} and \textcite{fourcade2009}.
Self-fulfilling beliefs are modeled in \textcite{schelling1978}.
On the crowding-out effect of incentives, see \textcite{frey1997}.

On the alternative cooperative anthropology, see \textcite{gintis2005} and \textcite{nowak2006}.
The behavioral economics challenge is surveyed in \textcite{kahneman2011}.
For a popular account of the ``selfish gene'' misreading, see Midgley (1979).

On cross-cultural variation in economic behavior, see the comprehensive survey in \textcite{henrich2001} and Henrich et al.\ (2010).
The neoliberal transformation of Japanese economic policy is discussed in Vogel (2006).
Nordic social democracy is analyzed in Esping-Andersen (1990).

\section{Exercises}

\begin{exercise}
Using Proposition~\ref{prop:funding-threshold}, estimate the annual counter-doxonomic investment needed to shift the U.S.\ anthropological regime from selfish to cooperative.
State your parameter assumptions explicitly.
Conduct a sensitivity analysis: how does the required investment change if $\alpha_C$ is twice as large (the cooperative message is more persuasive per dollar)?
\end{exercise}

\begin{exercise}
Design an experiment to test the self-fulfilling prophecy of selfishness narratives.
\begin{enumerate}[nosep]
  \item Randomly expose groups to different anthropological narratives.
  \item Measure cooperation in subsequent public goods games.
  \item Measure perceived norms (what do subjects think others will do?).
  \item Test whether the effect operates through the norm perception channel (mediation analysis).
\end{enumerate}
Specify sample size, treatment arms, and pre-registration plan.
\end{exercise}

\begin{exercise}
Construct a doxonomic account (Chapter~\ref{ch:belief-flow-accounting}) for the belief ``human nature is selfish,'' identifying at least five funding channels and estimating their annual expenditure.
Use the accounting framework to identify which channels have the highest return on doxonomic investment.
\end{exercise}

\begin{exercise}
Argue that the self-fulfilling property of selfishness beliefs creates a doxonomic trap (Definition~\ref{def:doxonomic-trap}).
\begin{enumerate}[nosep]
  \item Show formally that both the high-cooperation and low-cooperation equilibria are stable.
  \item Compute the basin of attraction for each equilibrium using the Beta(3,2) threshold distribution.
  \item Calculate the minimum doxonomic shock $\delta^*$ needed to tip the system from the low-cooperation to the high-cooperation equilibrium.
  \item Discuss what real-world intervention could produce a shock of this magnitude.
\end{enumerate}
\end{exercise}

\begin{exercise}
Compare the doxonomic investment in ``humans are selfish'' with the investment in ``humans are cooperative.''
\begin{enumerate}[nosep]
  \item Estimate the funding ratio using the institutional accounting approach.
  \item Given this ratio, what does the epistemic status of the dominant regime's dominance tell us---that it reflects the evidence, or that it reflects the funding?
  \item Design a study that could distinguish between these two explanations.
\end{enumerate}
\end{exercise}

\begin{exercise}
Conduct a content analysis of the first three chapters of two introductory economics textbooks.
\begin{enumerate}[nosep]
  \item Count explicit and implicit invocations of self-interest as a motivational assumption.
  \item Count mentions of cooperation, altruism, or pro-social behavior.
  \item Compare with the actual experimental evidence summarized in Section~\ref{sec:ch21-evidence}.
  \item What is the ``textbook-evidence gap,'' and how would you explain it doxonomically?
\end{enumerate}
\end{exercise}

\begin{exercise}
Select a non-Western society with a different anthropological regime (e.g., Japan, Bhutan, an indigenous commons community).
\begin{enumerate}[nosep]
  \item Describe the dominant anthropological assumptions.
  \item Map the institutional infrastructure that supports these assumptions (education, media, policy, religion).
  \item Assess whether neoliberal doxonomic intervention has shifted the regime, and if so, through which channels.
  \item What would the doxonomic model predict about the future trajectory?
\end{enumerate}
\end{exercise}

\begin{exercise}[Computational]
Simulate the self-fulfilling prophecy model from Theorem~\ref{thm:self-fulfilling}:
\begin{enumerate}[nosep]
  \item Generate 1,000 agents with thresholds $\tau_j \sim \text{Beta}(3, 2)$.
  \item Start at the high-cooperation equilibrium ($\hat{\mu} = 0.75$).
  \item Apply a doxonomic shock $\delta = 0.15$ and iterate the norm-updating dynamics for 50 periods.
  \item Plot the cooperation rate over time. Does the system tip to the low-cooperation equilibrium?
  \item Repeat for $\delta = 0.05, 0.10, 0.15, 0.20, 0.25$. Find the critical $\delta^*$ at which tipping occurs.
  \item Now simulate a ``counter-doxonomic'' intervention: starting from the low-cooperation equilibrium, apply a positive shock. What magnitude is required to restore high cooperation?
\end{enumerate}
\end{exercise}
